% from ega0-22-fr.tex
\oldpage[0\textsubscript{IV-1}]{195}

  \item[(ii)] Les éléments $d_Af_i\otimes1_{K'}$ forment une base du noyau
  $N=\Upsilon^{K'}_{A'/A}$ de l'homomorphisme canonique
  \[
    \Omega^1_A\otimes_AK'\longrightarrow
    \Omega^1_{A'}\otimes_{A'}K'.
  \]
\end{enumerate}
\end{proposition}

\begin{proof}
(i) Cette fois $A'\otimes_AK=A'/\mathfrak mA'$ est isomorphe à
$K[T_1,\ldots,T_r]/\mathfrak b'$, où $\mathfrak b'$ est l'idéal engendré
par les polynômes $T_i^p-a_i$~; l'hypothèse faite sur les $a_i$ entraîne que
cet anneau quotient est un corps (21.1.8), d'où les assertions de (i).

Avant de démontrer (ii), nous établirons le

\begin{lemma}[22.4.4.2]
\label{0.22.4.4.2-fr}
Soient $\Lambda$ un anneau, $A$ une $\Lambda$-algèbre, $\mathfrak J$ un
idéal de $A$ tel que l'anneau $C=A/\mathfrak J$ soit de caractéristique
$p>0$ (ce qui équivaut à dire que $p\cdot1\in\mathfrak J$ dans $A$). Soit
$\pi$ l'image canonique de $p\cdot1$ dans
$\mathfrak J/\mathfrak J^2$. Si $C$ est une $(\Lambda/p\Lambda)$-algèbre
formellement lisse pour les topologies discrètes, on a une suite exacte
canonique
\begin{equation}
\label{0.22.4.4.3-fr}
  0\longrightarrow(\mathfrak J/\mathfrak J^2)/C\cdot\pi
  \longrightarrow\Omega^1_{A/\Lambda}\otimes_AC
  \longrightarrow\Omega^1_{C/\Lambda}\longrightarrow0.
  \tag{22.4.4.3}
\end{equation}
\end{lemma}

\begin{proof}[Démonstration du lemme]
Posons en effet $A'=A/pA$, $\mathfrak J'=\mathfrak J/pA$, de sorte que
$C=A'/\mathfrak J'$. Comme
$A'=A\otimes_\Lambda(\Lambda/p\Lambda)$, on a
$\Omega^1_{A/\Lambda}\otimes_AA'=\Omega^1_{A'/\Lambda'}$ à un
isomorphisme canonique près, $\Lambda'$ désignant l'anneau
$\Lambda/p\Lambda$ (20.5.5). Appliquant alors la suite exacte (20.5.14.1) à
la $\Lambda'$-algèbre formellement lisse $C=A'/\mathfrak J'$, on obtient la
suite exacte
\[
  0\longrightarrow\mathfrak J'/\mathfrak J'^2
  \longrightarrow\Omega^1_{A'/\Lambda'}\otimes_{A'}C
  \longrightarrow\Omega^1_{C/\Lambda'}\longrightarrow0
\]
et comme l'homomorphisme $\Lambda\to\Lambda'$ est surjectif, on a
$\Omega^1_{C/\Lambda'}=\Omega^1_{C/\Lambda}$ à un isomorphisme canonique
près (20.5.15)~; d'où la suite exacte (22.4.4.3) en vertu de ce qui précède.
\end{proof}

(ii) Comme on a $\mathfrak m'^2=0$, nous désignerons encore par $V'$ l'idéal
$\mathfrak m'$ considéré comme $K'$-espace vectoriel, et nous poserons
\begin{equation}
\label{0.22.4.4.4-fr}
  V_0=V/K\cdot p=\mathfrak m/(\mathfrak m^2+pA),\qquad
  V_0'=V'/K'\cdot p=\mathfrak m'/(\mathfrak m'^2+pA').
  \tag{22.4.4.4}
\end{equation}
Appliquant la suite exacte (22.4.4.3) au cas où $\Lambda=\mathbf Z$, à la
$\mathbf Z$-algèbre $A$ (resp. $A'$) et à son idéal $\mathfrak m$
(resp. $\mathfrak m'$), il vient, puisque $K$ (resp. $K'$) est séparable sur
le corps premier $\mathbf Z/p\mathbf Z$, donc une
$(\mathbf Z/p\mathbf Z)$-algèbre formellement lisse (19.6.1), les suites
exactes
\begin{equation}
\label{0.22.4.4.5-fr}
  0\longrightarrow V_0\longrightarrow\Omega^1_A\otimes_AK
  \longrightarrow\Omega^1_K\longrightarrow0,\qquad
  0\longrightarrow V_0'\longrightarrow\Omega^1_{A'}\otimes_{A'}K'
  \longrightarrow\Omega^1_{K'}\longrightarrow0.
  \tag{22.4.4.5}
\end{equation}
Considérons alors le diagramme (22.4.4.6) où les deux colonnes médianes sont
la seconde suite (22.4.4.5) et la première tensorisée par $K'$~; comme la
première suite exacte (22.4.4.5) est formée d'espaces vectoriels sur $K$,
les deux colonnes médianes de (22.4.4.6) sont exactes~; en outre, le
diagramme formé par ces colonnes et les flèches horizontales qui les
connectent est commutatif, en vertu de la démonstration de (22.4.4.2) et de
la commutativité du diagramme (20.5.11.3). La dernière ligne de (22.4.4.6)
est la suite exacte obtenue à partir de (20.6.1.1) en y remplaçant $A$, $B$,
$C$ par $\mathbf Z$, $K$ et $K'$~; la ligne médiane est obtenue par
tensorisation avec $K'$ de la suite exacte de $K$-modules déduite de
(20.6.1.1) en y remplaçant $A$, $B$, $C$ par $\mathbf Z$, $A$, $K$
respectivement. Le diagramme formé de ces deux lignes et des flèches
verticales qui les connectent est commutatif en vertu de la commutativité de
(20.6.4.3). Notons d'autre part qu'on est sous les conditions d'application
de (22.4.1, (ii)), donc la ligne supérieure du diagramme (22.4.4.6) est
exacte, les colonnes extrêmes du diagramme

